\section{TreemapMix Augmentation Details}\label{app:labelmix-details}

This appendix provides implementation details omitted from the main Method section. We use the notation introduced in Section~\ref{sec:method}.

\begin{algorithm}[t]
\caption{TreemapMix augmentation}
\label{alg:labelmix}
\small
\begin{algorithmic}[1]
\Require dataset \(\mathcal{D}\), group size \(K\), concentration \(\alpha\), image size \(H\times W\)
\State build class-balanced rows and split them into groups of size \(K\)
\For{each group \(\mathcal{G}=\{(x_k,y_k)\}_{k=1}^{K}\)}
    \State sample \(w\sim\operatorname{Dirichlet}(\alpha,\ldots,\alpha)\)
    \State sort weights into \(\lambda_1\le\cdots\le\lambda_K\)
    \State construct boxes \(B_1,\dots,B_K\) with \(|B_i|/(HW)\approx\lambda_i\)
    \State apply a random valid symmetry transform to the boxes
    \For{\(s=0,\ldots,K-1\)}
        \State assign each source image \(x_{\pi_s(i)}\) to slot \(i\), for \(i=1,\dots,K\)
        \State render \(\tilde{x}^{(s)}=\sum_i M_i\odot\operatorname{Resize}(x_{\pi_s(i)},B_i)\)
        \State output \(\bigl(\tilde{x}^{(s)},\{(y_{\pi_s(i)},\lambda_i)\}_{i=1}^{K}\bigr)\)
    \EndFor
\EndFor
\end{algorithmic}
\end{algorithm}

\subsection{Class-balanced row construction}

Let the training dataset be \(\mathcal{D}=\{(x_n,y_n)\}_{n=1}^{N_{\mathrm{data}}}\), where \(y_n\in\{1,\dots,C\}\), and let \(\mathcal{I}_c=\{n:y_n=c\}\) be the set of dataset indices belonging to class $c$. A class-balanced row is constructed by sampling one example from each class. For imbalanced datasets, minority classes may be sampled with replacement, while majority classes may be subsampled across rows, so that every class contributes equally to each row.

Each row is shuffled and split into groups of size $K$. When $K\le C$, complete groups formed from a class-balanced row contain distinct labels by construction, because each class appears at most once in the row. This encourages diverse multi-class mixtures and avoids repeated target classes in most TreemapMix groups.

\subsection{Incomplete groups and row-size preservation}

Suppose a row contains $N_{\mathrm{row}}$ samples. Write \(N_{\mathrm{row}}=K\left\lfloor\frac{N_{\mathrm{row}}}{K}\right\rfloor+r\), where \(0\le r<K\). TreemapMix processes the first \(K\left\lfloor\frac{N_{\mathrm{row}}}{K}\right\rfloor\)
samples as complete groups of size $K$. If $r>0$, the remaining $r$ samples are padded with $K-r$ samples drawn without replacement from earlier entries in the same row, and only the first $r$ cyclic outputs are retained. This padding can repeat labels. In natural-frequency mode, rows shorter than $K$ are passed through as unmixed samples. The padding can repeat labels: SCE sums their contributions, while the PL implementation processes them as separate entries without aggregation.

TreemapMix generates $K$ cyclic views from the temporary group but retains only the first $r$ outputs. Therefore, the total number of final TreemapMix samples contributed by the row is \(K\left\lfloor\frac{N_{\mathrm{row}}}{K}\right\rfloor+r=N_{\mathrm{row}}\).
Thus, the number of TreemapMix outputs matches the number of samples in the constructed row stream, even when a row size is not divisible by $K$.

\subsection{Integer box conversion and layout fallback}

The squarified treemap is first computed in normalized coordinates. Let \(\bar{B}_i=(\bar{x}_i,\bar{y}_i,\bar{w}_i,\bar{h}_i)\) denote the normalized rectangle for slot $i$. We convert it to an integer pixel box \(B_i=(x_i^{(0)},y_i^{(0)},x_i^{(1)},y_i^{(1)})\). Because independent rounding can introduce gaps or overlaps, shared boundaries are snapped consistently. After conversion, the boxes satisfy \(B_i\cap B_j=\emptyset\) for \(i\neq j\) and \(\bigcup_{i=1}^{K}B_i=[0,W]\times[0,H]\).

A converted layout is considered valid if every box has positive area. Equivalently, for integer boxes, every box must have width and height at least one pixel. If the treemap layout fails this validity check, TreemapMix falls back to a stripe layout. The stripe layout partitions the canvas along one axis according to the sorted weights $\lambda_1,\dots,\lambda_K$, then applies the same integer-boundary snapping procedure. This fallback preserves a valid full-canvas tiling and maintains the sampled area weights up to rounding error.

\subsection{Symmetry transformations}

For square images, TreemapMix samples from the dihedral group
\[
\begin{aligned}
D_4=\{&
\text{identity},\space
90^\circ\text{ rotation},\space
180^\circ\text{ rotation},\space
270^\circ\text{ rotation},\\
&
\text{horizontal flip},\space
\text{vertical flip},\space
\text{transpose},\space
\text{anti-transpose}
\}.
\end{aligned}
\]
For non-square images, TreemapMix uses the shape-preserving subset
\[
\{
\text{identity},\space
180^\circ\text{ rotation},\space
\text{horizontal flip},\space
\text{vertical flip}
\}.
\]
These transformations preserve region areas, so they increase spatial diversity without changing the TreemapMix target.

\subsection{Cyclic views and slot exposure}\label{app:cyclic-views}

Cyclic views are a tractable alternative to enumerating all $K!$ image-to-slot assignments. If an epoch contains $G$ complete TreemapMix groups, full permutation enumeration would produce $K!G$ mixed images, increasing the sample count by \(\frac{K!G}{KG}=(K-1)!\). In contrast, cyclic TreemapMix produces exactly $K$ views per group, giving $KG$ mixed images and preserving the number of training samples after class-balanced row construction.

Cyclic permutation is designed to decouple source labels from fixed spatial slots. Although any individual mixed view may assign a class to a small or large region, the full cyclic orbit exposes every source image to every slot exactly once. Therefore, when the source labels in a group are distinct, every participating class is paired once with every area weight and, for PL training, every rank position induced by the sorted slot weights. Appendix Lemma~\ref{lem:cyclic-balancing} formalizes the resulting slot-exposure balance.

\section{Soft TreemapMix Cross-Entropy Derivations}\label{app:soft-ce}

We use the sparse target notation from Section~\ref{sec:soft-ce}. For one TreemapMix sample, let \(\mathcal{T}=\{(c_i,\eta_i)\}_{i=1}^{K}\), where \(\eta_i\ge 0\) and \(\sum_{i=1}^{K}\eta_i=1\),
where $c_i$ is the class assigned to slot $i$ and $\eta_i$ is the corresponding slot weight. For cyclic view $s$, we have $c_i=y_{\pi_s(i)}$ and $\eta_i=\lambda_i$.

\subsection{KL derivation}

The sparse target distribution is \(q_c=\sum_{i=1}^{K}\eta_i\mathbf{1}[c_i=c]\).
It is a valid probability distribution because
\[
\sum_{c=1}^{C}q_c
=
\sum_{c=1}^{C}
\sum_{i=1}^{K}
\eta_i\mathbf{1}[c_i=c]
=
\sum_{i=1}^{K}
\eta_i
\sum_{c=1}^{C}
\mathbf{1}[c_i=c]
=
\sum_{i=1}^{K}\eta_i
=
1.
\]
Hence $q=(q_1,\dots,q_C)$ lies on the probability simplex.

The KL divergence from $q$ to $p_\theta$ is
\[
D_{\mathrm{KL}}(q\|p_\theta)
=
\sum_{c=1}^{C}
q_c\log\frac{q_c}{p_\theta(c\mid\tilde{x})}
=
\sum_{c=1}^{C}q_c\log q_c
-
\sum_{c=1}^{C}q_c\log p_\theta(c\mid\tilde{x}).
\]
The first term is independent of $\theta$, so minimizing the KL divergence is equivalent to minimizing the dense soft cross-entropy
\[
\mathcal{L}_{\mathrm{TM\text{-}CE}}
=
-\sum_{c=1}^{C}
q_c\log p_\theta(c\mid\tilde{x}).
\]
Substituting the sparse target definition gives
\[
\begin{aligned}
\mathcal{L}_{\mathrm{TM\text{-}CE}}
&=
-\sum_{c=1}^{C}
\left(
\sum_{i=1}^{K}
\eta_i\mathbf{1}[c_i=c]
\right)
\log p_\theta(c\mid\tilde{x}) \\
&=
-\sum_{i=1}^{K}
\eta_i\log p_\theta(c_i\mid\tilde{x}),
\end{aligned}
\]
which is the sparse form in \eqref{eq:lm-ce}. For a mini-batch of size $B$,
\[
\mathcal{L}_{\mathrm{TM\text{-}CE}}
=
-\frac{1}{B}
\sum_{b=1}^{B}
\sum_{i=1}^{K}
\eta_{b,i}\log p_\theta(c_{b,i}\mid\tilde{x}_b).
\]

\subsection{Gradient derivation}

For one TreemapMix sample, write
\[
\mathcal{L}_{\mathrm{TM\text{-}CE}}
=
-\sum_{c=1}^{C}q_c\log p_c,
\qquad
p_c=
\frac{\exp(z_c)}{\sum_{j=1}^{C}\exp(z_j)}.
\]
Using
\[
\log p_c
=
z_c-\log\sum_{j=1}^{C}\exp(z_j),
\]
we obtain
\[
\mathcal{L}_{\mathrm{TM\text{-}CE}}
=
-\sum_{c=1}^{C}q_c z_c
+
\log\sum_{j=1}^{C}\exp(z_j),
\]
because $\sum_{c=1}^{C}q_c=1$. Differentiating with respect to $z_i$ gives
\[
\frac{\partial \mathcal{L}_{\mathrm{TM\text{-}CE}}}{\partial z_i}
=
-q_i
+
\frac{\exp(z_i)}{\sum_{j=1}^{C}\exp(z_j)}
=
p_i-q_i.
\]
Using the sparse target definition,
\[
\frac{\partial \mathcal{L}_{\mathrm{TM\text{-}CE}}}{\partial z_i}
=
p_\theta(i\mid\tilde{x})
-
\sum_{j=1}^{K}
\eta_j\mathbf{1}[c_j=i].
\]
For a mini-batch of size $B$,
\[
\frac{\partial \mathcal{L}_{\mathrm{TM\text{-}CE}}}{\partial z_{b,c}}
=
\frac{1}{B}
\left[
p_\theta(c\mid\tilde{x}_b)
-
\sum_{i=1}^{K}
\eta_{b,i}\mathbf{1}[c_{b,i}=c]
\right].
\]

\subsection{Random cyclic-view target variance}

\refstepcounter{lemma}\label{lem:ce-cyclic-variance}
\paragraph{Lemma~\thelemma\ (Random cyclic-view and cyclic-average target variance).}
Assume the $K$ source labels in a TreemapMix group are distinct, and let \(q_c^{(s)}=\sum_{i=1}^{K}\lambda_i\mathbf{1}[y_{\pi_s(i)}=c]\) be the TreemapMix-SCE target mass for class $c$ in cyclic view $s$. Conditioned on the source labels, if $c$ is a participating class and $s$ is sampled uniformly from $\{0,\dots,K-1\}$, then \(\mathbb{E}_{w,s}[q_c^{(s)}]=\frac{1}{K}\) and \(\operatorname{Var}_{w,s}(q_c^{(s)})=\frac{K-1}{K^2(K\alpha+1)}\). Moreover, averaging over the full cyclic orbit gives \(\frac{1}{K}\sum_{s=0}^{K-1}q_c^{(s)}=\frac{1}{K}\), so the cyclic-average target mass has zero variance conditioned on the source labels. For an absent class, \(q_c^{(s)}=0\) for all \(s\).

\paragraph{Proof.}

We prove Lemma~\ref{lem:ce-cyclic-variance}. Assume the $K$ source labels are distinct, and let $c=y_k$ be a participating class. For each cyclic view $s$, class $c$ appears in exactly one slot. Since the cyclic map visits every slot exactly once as $s$ ranges from $0$ to $K-1$, sampling $s$ uniformly makes the slot index containing class $c$ uniformly distributed over $\{1,\dots,K\}$. Therefore, \(q_c^{(s)}=\lambda_I\), where \(I\sim\operatorname{Uniform}\{1,\dots,K\}\).
Because $(\lambda_1,\dots,\lambda_K)$ is a sorted permutation of $(w_1,\dots,w_K)$, a uniformly selected sorted component $\lambda_I$ has the same marginal distribution as a uniformly selected unsorted component $w_J$. By exchangeability of the symmetric Dirichlet distribution, this marginal distribution is the same as that of any component $w_j$. Hence
\(\mathbb{E}_{w,s}[q_c^{(s)}]=\mathbb{E}[w_j]=\frac{1}{K}\).
For a symmetric Dirichlet distribution with concentration parameter $\alpha$ in each coordinate,
\(\operatorname{Var}(w_j)=\frac{K-1}{K^2(K\alpha+1)}\). Therefore, \(\operatorname{Var}_{w,s}(q_c^{(s)})=\operatorname{Var}(w_j)=\frac{K-1}{K^2(K\alpha+1)}\).

The cyclic-average identity follows from Lemma~\ref{lem:cyclic-balancing}. When labels are distinct, \(m_c=1\) for every participating class, so \(\frac{1}{K}\sum_{s=0}^{K-1}q_c^{(s)}=\frac{1}{K}\).
Since this equality holds for every Dirichlet sample, the conditional variance of the cyclic-average target mass is zero. If $c$ is absent from the group, then no slot carries class $c$, and \(q_c^{(s)}=0\) for all \(s\).

Finally, if logits are fixed, then \(\frac{\partial \mathcal{L}_{\mathrm{TM\text{-}CE}}}{\partial z_c}=p_c-q_c^{(s)}\)
has the same variance as \(q_c^{(s)}\), since \(p_c\) is fixed under the conditioning. This proves that a randomly selected cyclic view preserves Dirichlet-controlled target variability, while the full cyclic orbit exactly balances average target mass.

\section{TreemapMix Plackett--Luce Derivations}\label{app:pl}

The derivation assumes $K$ distinct target classes. Let $r_1\succ\cdots\succ r_K$ denote their descending order by area weight, set $L=K$, and write $\rho_{r_\ell}$ for the area weight assigned to class $r_\ell$. Let $\mathcal{U}=\{1,\dots,C\}\setminus\{r_1,\dots,r_K\}$ be the absent classes, with total predicted probability mass $p_{\mathcal U}=1-\sum_{j=1}^{K}p_{r_j}$.

\subsection{From the Plackett--Luce probability to TreemapMix-PL}

Let the model output logits \(z=f_\theta(\tilde{x})\in\mathbb{R}^{C}\) and softmax probabilities \(p_c=p_\theta(c\mid\tilde{x})\).
For the descending ranking $r_1\succ\cdots\succ r_L$, the top-$L$ Plackett--Luce probability with unranked residual mass is
\[
P_{\mathrm{PL}}(r_1\succ\cdots\succ r_L\mid\tilde{x})
=
\prod_{\ell=1}^{L}
\frac{
p_{r_\ell}
}{
p_{\mathcal U}
+
\sum_{j=\ell}^{L}p_{r_j}
}.
\]
The corresponding unweighted ListMLE-style negative log-likelihood is
\[
\mathcal{L}_{\mathrm{ListMLE}}
=
-\sum_{\ell=1}^{L}
\log
\frac{
p_{r_\ell}
}{
p_{\mathcal U}
+
\sum_{j=\ell}^{L}p_{r_j}
}.
\]
TreemapMix-PL weights each ranking term by its area weight $\eta_{r_\ell}$:
\[
\mathcal{L}_{\mathrm{TM\text{-}PL}}
=
-\sum_{\ell=1}^{L}
\rho_{r_\ell}
\log
\frac{
p_{r_\ell}
}{
p_{\mathcal U}
+
\sum_{j=\ell}^{L}p_{r_j}
}.
\]
This is the descending form in \eqref{eq:lm-pl-desc}. When all target weights are equal, this objective is a constant rescaling of the unweighted ListMLE loss.

\subsection{Ascending implementation form}

For implementation, we reverse the descending order. Define \(a_i=r_{L+1-i}\), \(\gamma_i=\rho_{a_i}\), and \(\gamma_1\le\gamma_2\le\cdots\le\gamma_L\).
Thus, $a_1$ is the weakest ranked class and $a_L$ is the strongest. Reindexing the descending objective gives
\[
\mathcal{L}_{\mathrm{TM\text{-}PL}}
=
-\sum_{i=1}^{L}
\gamma_i
\log
\frac{
p_{a_i}
}{
p_{\mathcal U}
+
\sum_{j=1}^{i}p_{a_j}
}.
\]
Using
\[
D_i
=
p_{\mathcal U}
+
\sum_{j=1}^{i}p_{a_j}
=
1-\sum_{j=i+1}^{L}p_{a_j},
\]
we obtain
\[
\mathcal{L}_{\mathrm{TM\text{-}PL}}
=
\sum_{i=1}^{L}
\gamma_i
\left[
\log D_i-\log p_{a_i}
\right],
\]
which is the ascending implementation form in \eqref{eq:lm-pl-asc}. For a mini-batch of size $B$, where sample $b$ has $L_b$ valid ranked entries, the implementation uses
\[
\mathcal{L}_{\mathrm{TM\text{-}PL}}
=
\frac{1}{B}
\sum_{b=1}^{B}
\sum_{i=1}^{L_b}
\gamma_{b,i}
\left[
\log D_{b,i}
-
\log p_\theta(a_{b,i}\mid \tilde{x}_b)
\right].
\]

\subsection{Gradient derivation}

For one sample,
\[
\mathcal{L}_{\mathrm{TM\text{-}PL}}
=
\sum_{i=1}^{L}
\gamma_i
\left[
\log D_i-\log p_{a_i}
\right].
\]
Define the stronger-class probability mass \(R_i=\sum_{j=i+1}^{L}p_{a_j}\) and \(D_i=1-R_i\), and the stronger-class set \(T_i=\{a_{i+1},\dots,a_L\}\). First, \(\frac{\partial \log p_{a_i}}{\partial z_c}=\mathbf{1}[c=a_i]-p_c\).
Next,
\[
\begin{aligned}
\frac{\partial R_i}{\partial z_c}
&=
\sum_{j=i+1}^{L}
p_{a_j}
\left(
\mathbf{1}[c=a_j]-p_c
\right) \\
&=
p_c\mathbf{1}[c\in T_i]-p_cR_i \\
&=
p_c\left(\mathbf{1}[c\in T_i]-R_i\right).
\end{aligned}
\]
Therefore,
\[
\frac{\partial D_i}{\partial z_c}
=
-\frac{\partial R_i}{\partial z_c}
=
p_c\left(R_i-\mathbf{1}[c\in T_i]\right),
\]
and hence
\[
\frac{\partial \log D_i}{\partial z_c}
=
\frac{
p_c\left(R_i-\mathbf{1}[c\in T_i]\right)
}{
D_i
}.
\]
For the $i$-th loss term
\[
\ell_i
=
\gamma_i
\left[
\log D_i-\log p_{a_i}
\right],
\]
we have
\[
\begin{aligned}
\frac{\partial \ell_i}{\partial z_c}
&=
\gamma_i
\left[
\frac{
p_c\left(R_i-\mathbf{1}[c\in T_i]\right)
}{
D_i
}
-
\mathbf{1}[c=a_i]
+
p_c
\right] \\
&=
\gamma_i
\left[
\frac{
p_c\left(R_i-\mathbf{1}[c\in T_i]+D_i\right)
}{
D_i
}
-
\mathbf{1}[c=a_i]
\right] \\
&=
\gamma_i
\left[
\frac{
p_c\left(1-\mathbf{1}[c\in T_i]\right)
}{
D_i
}
-
\mathbf{1}[c=a_i]
\right],
\end{aligned}
\]
because $R_i+D_i=1$. Summing over $i$ gives
\[
\frac{\partial \mathcal{L}_{\mathrm{TM\text{-}PL}}}{\partial z_c}
=
\sum_{i=1}^{L}
\gamma_i
\left[
\frac{
p_c\left(1-\mathbf{1}[c\in T_i]\right)
}{
D_i
}
-
\mathbf{1}[c=a_i]
\right],
\]
which is the gradient in \eqref{eq:lm-pl-grad}.

\subsection{Ordinal gradient structure}

The gradient above has a simple ordinal interpretation.

\subsubsection*{Unranked class}

Let $u\notin\{a_1,\dots,a_L\}$. Then \(\mathbf{1}[u\in T_i]=0\) and \(\mathbf{1}[u=a_i]=0\) for all $i$, so \(\frac{\partial \mathcal{L}_{\mathrm{TM\text{-}PL}}}{\partial z_u}=p_u\sum_{i=1}^{L}\frac{\gamma_i}{D_i}\ge 0\).
Thus, gradient descent suppresses classes that do not appear in the TreemapMix target.

\subsubsection*{Ranked class}

Let $c=a_m$ for some $m\in\{1,\dots,L\}$. For $i<m$, class $a_m$ belongs to the stronger set $T_i$, so those terms vanish. For $i=m$, the contribution is \(\gamma_m\left(\frac{p_{a_m}}{D_m}-1\right)\). For $i>m$, class $a_m$ is not in $T_i$ and is not equal to $a_i$, so the remaining contribution is \(p_{a_m}\sum_{i=m+1}^{L}\frac{\gamma_i}{D_i}\).
Therefore,
\[
\frac{\partial \mathcal{L}_{\mathrm{TM\text{-}PL}}}{\partial z_{a_m}}
=
\gamma_m
\left(
\frac{p_{a_m}}{D_m}-1
\right)
+
p_{a_m}
\sum_{i=m+1}^{L}
\frac{\gamma_i}{D_i}.
\]
For the strongest class $a_L$, the sum is empty and $D_L=1$, so \(\frac{\partial \mathcal{L}_{\mathrm{TM\text{-}PL}}}{\partial z_{a_L}}=\gamma_L(p_{a_L}-1)\le 0\).
Thus, gradient descent promotes the strongest ranked class unless $p_{a_L}=1$, suppresses unranked classes, and regulates intermediate ranked classes through competition with stronger classes.

\subsection{Proof of the Dirichlet order-statistics result}

\refstepcounter{lemma}\label{lem:pl-order}
\paragraph{Lemma~\thelemma\ (Dirichlet order statistics and ranking emphasis).}
Assume the $K$ target labels are distinct, so that $L=K$ and $\gamma_i=\lambda_i$. Let \(\mu_i(\alpha)=\mathbb{E}[\lambda_i]\) denote the expected $i$-th sorted Dirichlet weight. Conditioned on fixed logits and fixed rank assignments, the expected TreemapMix-PL gradient for a ranked class $a_m$ satisfies \(\mathbb{E}_{w}\left[\frac{\partial \mathcal{L}_{\mathrm{TM\text{-}PL}}}{\partial z_{a_m}}\right]=\mu_m(\alpha)\left(\frac{p_{a_m}}{D_m}-1\right)+p_{a_m}\sum_{i=m+1}^{K}\frac{\mu_i(\alpha)}{D_i}\), and for an unranked class $u$, \(\mathbb{E}_{w}\left[\frac{\partial \mathcal{L}_{\mathrm{TM\text{-}PL}}}{\partial z_u}\right]=p_u\sum_{i=1}^{K}\frac{\mu_i(\alpha)}{D_i}\).

\paragraph{Proof.}

Assume the $K$ target labels are distinct, so that $L=K$ and the ascending PL weights satisfy \(\gamma_i=\lambda_i\). Let \(\mu_i(\alpha)=\mathbb{E}[\lambda_i]\)
denote the expected $i$-th sorted Dirichlet weight. Conditioned on fixed logits and fixed rank assignments, $D_i$ and $p_c$ are fixed, and the ranked-class gradient is linear in $\lambda_i$. Therefore,
\[
\mathbb{E}_{w}
\left[
\frac{\partial \mathcal{L}_{\mathrm{TM\text{-}PL}}}{\partial z_{a_m}}
\right]
=
\mu_m(\alpha)
\left(
\frac{p_{a_m}}{D_m}-1
\right)
+
p_{a_m}
\sum_{i=m+1}^{K}
\frac{\mu_i(\alpha)}{D_i}.
\]
Likewise, for an unranked class $u$,
\[
\mathbb{E}_{w}
\left[
\frac{\partial \mathcal{L}_{\mathrm{TM\text{-}PL}}}{\partial z_u}
\right]
=
p_u
\sum_{i=1}^{K}
\frac{\mu_i(\alpha)}{D_i}.
\]
Thus, in the distinct-label case, $\alpha$ affects TreemapMix-PL through the order statistics of the sampled Dirichlet weights. Smaller $\alpha$ widens the expected gap between low- and high-rank weights, concentrating supervision on dominant labels. As $\alpha$ increases, the sorted weights concentrate near $1/K$, so the loss approaches a constant rescaling of an unweighted ListMLE-style objective. This proves that, in the distinct-label case, the expected TreemapMix-PL gradient depends on $\alpha$ through the sorted Dirichlet weights assigned to rank positions.

\section{Proof of Cyclic Permutation Balancing}\label{app:cyclic-balancing}

\refstepcounter{lemma}\label{lem:cyclic-balancing}
\paragraph{Lemma~\thelemma\ (Cyclic permutation balances slot and target exposure)}

Let \refstepcounter{equation}\label{eq:cyclic-soft-target}\(q_c^{(s)}=\sum_{i=1}^{K}\lambda_i \mathbf{1}[y_{\pi_s(i)}=c], m_c=\sum_{k=1}^{K}\mathbf{1}[y_k=c]\)\textup{ (\theequation)}. Then, for every class $c$ and slot $i$, \refstepcounter{equation}\label{eq:slot-balance}\(\frac{1}{K}\sum_{s=0}^{K-1}\mathbf{1}[y_{\pi_s(i)}=c]=\frac{m_c}{K}\)\textup{ (\theequation)}, and consequently \refstepcounter{equation}\label{eq:cyclic-balance-soft}\(\frac{1}{K}\sum_{s=0}^{K-1}q_c^{(s)}=\frac{m_c}{K}\)\textup{ (\theequation)}.

\paragraph{Proof.}

For cyclic view $s$, slot $i$ receives source image $x_{\pi_s(i)}$ and label $y_{\pi_s(i)}$, where \(\pi_s(i)=1+\bigl((i+s-1)\bmod K\bigr)\).
For any fixed slot $i$, as $s$ ranges from $0$ to $K-1$, the index $\pi_s(i)$ visits every source position exactly once. Therefore, for any class $c$,
\(\sum_{s=0}^{K-1}\mathbf{1}[y_{\pi_s(i)}=c]=\sum_{k=1}^{K}\mathbf{1}[y_k=c]=m_c\). Dividing by $K$ gives the slot-balancing identity \(\frac{1}{K}\sum_{s=0}^{K-1}\mathbf{1}[y_{\pi_s(i)}=c]=\frac{m_c}{K}\),
which proves \eqref{eq:slot-balance}.

Next, the soft target mass assigned to class $c$ in cyclic view $s$ is \(q_c^{(s)}=\sum_{i=1}^{K}\lambda_i\mathbf{1}[y_{\pi_s(i)}=c]\).
Averaging over cyclic views gives
\[
\begin{aligned}
\frac{1}{K}
\sum_{s=0}^{K-1}
q_c^{(s)}
&=
\frac{1}{K}
\sum_{s=0}^{K-1}
\sum_{i=1}^{K}
\lambda_i\mathbf{1}[y_{\pi_s(i)}=c] \\
&=
\frac{1}{K}
\sum_{i=1}^{K}
\lambda_i
\sum_{s=0}^{K-1}
\mathbf{1}[y_{\pi_s(i)}=c] \\
&=
\frac{1}{K}
\sum_{i=1}^{K}
\lambda_i m_c \\
&=
\frac{m_c}{K},
\end{aligned}
\]
because $\sum_{i=1}^{K}\lambda_i=1$. This proves \eqref{eq:cyclic-balance-soft}.

When the $K$ source labels are distinct, $m_c=1$ for every participating class and $m_c=0$ for every absent class. Therefore, each participating class appears exactly once in every slot across the $K$ cyclic views. Since slots are associated with sorted weights and rank positions, cyclic permutation also prevents any participating class from being permanently tied to a fixed TreemapMix weight or rank. The PL denominator terms remain prediction-dependent, so this balancing does not remove ordinal competition; it only removes persistent source-label-to-slot bias.

\section{Implementation Details}
\subsection{Experimental Protocol and Hyperparameters}
\label{app:appendix-aug-configs}

Tables~\ref{tab:timm-adaptations}, \ref{tab:baseline-hparams-aug}, and
\ref{tab:baseline-hparams-train} summarize the ImageNet-1K experimental
protocol and baseline hyperparameters. Table~\ref{tab:timm-adaptations}
summarizes the schedule-level adaptations applied to the reference
\texttt{timm} recipes for the main 110-epoch comparison.
Tables~\ref{tab:baseline-hparams-aug} and
\ref{tab:baseline-hparams-train} separate data-augmentation settings from
optimizer and training settings for readability.

The baseline recipes are adapted from established \texttt{timm} ImageNet-1K
training configurations~\citep{rw2019timm,262588213843476SearchingBetterVit}.
Because our main study uses a shorter 110-epoch budget than many reference
recipes, we adapt the training recipe: we disable
\texttt{aug\_repeats}, adjust the learning rate, disable gradient clipping, and keep the remaining regularization settings fixed across methods.
All methods therefore share the same architecture, data budget, optimizer schedule, and evaluation protocol. The comparison isolates augmentation and loss effects, but does not claim per-baseline optimality.

The main ImageNet-1K comparison uses a global batch size of 1024. We omit epochs and warmup epochs
from Tables~\ref{tab:baseline-hparams-aug} and
\ref{tab:baseline-hparams-train} because those values are held fixed across
runs. The augmentation table lists the ImageNet-1K augmentation recipes used
in the main comparisons.

For the main comparison, RICAP uses the following auxiliary settings: the augmentation is
applied with probability 0.5, disabled for the final 10 epochs, allowed to
vary the crop center over a ratio range from 0.5 to 1.5, and allowed to vary
the crop scale from 0.8 to 1.2. In this comparison, TreemapMix-SCE and
TreemapMix-mixed runs use $K=4$ with a cosine $\alpha$ schedule from $0.1$
to $0.5$, while TreemapMix-PL runs use $K=6$ with the same cosine
$\alpha$ schedule.

\begin{table}[t]
  \centering
  \setlength{\tabcolsep}{5pt}
  \caption{Adaptations to the reference \texttt{timm} ImageNet-1K recipes used
for the main 110-epoch comparison. We shorten the original long schedules,
disable augmentation repeats, switch to \texttt{bfloat16}, and retune the
learning rate; all other recipe settings are kept unchanged unless listed.}
  \label{tab:timm-adaptations}
  \begin{tabular}{llll}
    \toprule
    Model & Setting & Original & New \\
    \midrule
    \multicolumn{4}{l}{\textbf{ViT-Betwixt}} \\
    & \texttt{dtype} & \texttt{float16} & \texttt{bfloat16} \\
    & \texttt{aug\_repeats} & 3.0 & 0.0 \\
    & \texttt{warmup\_epochs} & 20 & 10 \\
    & \texttt{epochs} & 600 & 110 \\
    & \texttt{lr} & 1.2e-3 & 1.5e-3 \\
    \midrule
    \multicolumn{4}{l}{\textbf{ViT-Medium}} \\
    & \texttt{dtype} & \texttt{float16} & \texttt{bfloat16} \\
    & \texttt{aug\_repeats} & 3.0 & 0.0 \\
    & \texttt{clip\_grad} & 5.0 & \texttt{null} \\
    & \texttt{warmup\_epochs} & 20 & 10 \\
    & \texttt{epochs} & 600 & 110 \\
    & \texttt{lr} & 1.2e-3 & 1.5e-3 \\
    \midrule
    \multicolumn{4}{l}{\textbf{ViT-Little}} \\
    & \texttt{dtype} & \texttt{float16} & \texttt{bfloat16} \\
    & \texttt{aug\_repeats} & 3.0 & 0.0 \\
    & \texttt{clip\_grad} & 5.0 & \texttt{null} \\
    & \texttt{warmup\_epochs} & 20 & 10 \\
    & \texttt{epochs} & 600 & 110 \\
    & \texttt{lr} & 1.2e-3 & 1.5e-3 \\
    \midrule
    \multicolumn{4}{l}{\textbf{ViT-Wee}} \\
    & \texttt{dtype} & \texttt{float16} & \texttt{bfloat16} \\
    & \texttt{aug\_repeats} & 3.0 & 0.0 \\
    & \texttt{clip\_grad} & 5.0 & \texttt{null} \\
    & \texttt{warmup\_epochs} & 20 & 10 \\
    & \texttt{epochs} & 450 & 110 \\
    & \texttt{lr} & 1.2e-3 & 1.5e-3 \\
    \bottomrule
  \end{tabular}
\end{table}

\begin{table*}[t]
  \centering
  \caption{ViT data-augmentation hyperparameters for the ImageNet-1K main
  comparison recipes. A checkmark in the Single-Image Aug. column indicates
  that the recipe includes standard image-only transforms such as cropping,
  flipping, scaling, color perturbations, AutoAugment-style transformations
  \citep{cubukRandAugmentPracticalAutomated2019}, and random erasing. The
  fixed auxiliary RICAP and TreemapMix settings are described in
  the appendix text.}
  \label{tab:baseline-hparams-aug}

  \begin{tabular}{lccccccc}
    \toprule
    Config
      & \makecell{Single-Image\\Aug}
      & Mixup
      & CutMix
      & \makecell{Mixup\\Prob}
      & \makecell{CutMix\\Prob}
      & \makecell{RICAP\\Prob} \\
    \midrule
    No Aug
      & $\times$
      & 0.0
      & 0.0
      & 0.0
      & 0.0
      & 0.0 \\

    \makecell[l]{Single-Image Aug}
      & $\checkmark$
      & 0.0
      & 0.0
      & 0.0
      & 0.0
      & 0.0 \\

    Mixup
      & $\checkmark$
      & 0.8
      & 0.0
      & 1.0
      & 0.0
      & 0.0 \\

    CutMix
      & $\checkmark$
      & 0.0
      & 1.0
      & 0.0
      & 1.0
      & 0.0 \\

    \makecell[l]{Mixup+CutMix}
      & $\checkmark$
      & 0.8
      & 1.0
      & 0.5
      & 0.5
      & 0.0 \\

    RICAP
      & $\checkmark$
      & 0.0
      & 0.0
      & 0.0
      & 0.0
      & 0.5 \\

    \midrule
    \makecell[l]{TreemapMix-SCE}
      & $\checkmark$
      & 0.0
      & 0.0
      & 0.0
      & 0.0
      & 0.0 \\

    \makecell[l]{TreemapMix-PL}
      & $\checkmark$
      & 0.0
      & 0.0
      & 0.0
      & 0.0
      & 0.0 \\

    \makecell[l]{TreemapMix-mixed}
      & $\checkmark$
      & 0.0
      & 0.0
      & 0.0
      & 0.0
      & 0.0 \\
    \bottomrule
  \end{tabular}%
\end{table*}

\begin{table*}[!htbp]
  \centering
  \caption{ViT baseline optimizer and training hyperparameters for the main comparison. Opt. denotes optimizer, LR learning rate, WD
  weight decay, and EMA exponential moving average of model weights. NAdamW
  is Nesterov-accelerated Adam with decoupled weight decay. DropPath is the
  stochastic-depth probability.}
  \label{tab:baseline-hparams-train}
  \begin{tabular}{lcccccc}
    \toprule
    Model & Opt. & Base LR & LR Decay & WD & DropPath & EMA \\
    \midrule
    ViT-Wee & NAdamW & 1.5e-3 & cosine & 0.06 & 0.2 & yes \\
    ViT-Little & NAdamW & 1.5e-3 & cosine & 0.06 & 0.2 & no \\
    ViT-Medium & NAdamW & 1.5e-3 & cosine & 0.06 & 0.2 & no \\
    ViT-Betwixt & NAdamW & 1.5e-3 & cosine & 0.08 & 0.2 & no \\
    \bottomrule
  \end{tabular}%
\end{table*}

\subsection{COCO transfer setup}
\label{app:det-hparams}

Table~\ref{tab:det-hparams} summarizes the ViTDet-style configurations used
for the COCO transfer experiments. In all cases the starting point
is set to the corresponding checkpoint from the main classification
experiments, and we fine-tune a Mask R-CNN model with a ViTDet-style adapter
for 30 epochs on COCO.

\begin{table*}[!htbp]
  \centering
  \small
  \caption{Hyperparameters for the object-detection transfer experiments with
  a ViTDet-style adapter and Mask R-CNN on COCO. Embed Dim is the embedding
  dimension, Heads is the attention-head count, MLP Ratio is the multilayer
  perceptron expansion ratio, and LR is the learning rate.}
  \label{tab:det-hparams}
  \begin{tabular}{lccccccc}
    \toprule
    Model & Embed Dim & Depth & Heads & MLP Ratio & Drop Path & LR & Layer Decay \\
    \midrule
    ViT-Wee & 256 & 14 & 4 & 5.0 & 0.1 & 6.0e-4 & 0.7 \\
    ViT-Betwixt & 640 & 12 & 10 & 4.0 & 0.1 & 8.0e-4 & 0.7 \\
    \bottomrule
  \end{tabular}%
\end{table*}

\subsection{Loss functions in code}
\label{app:loss-code}

We include concise PyTorch reference implementations of TreemapMix-SCE
(Figure~\ref{fig:treemapmix-sce-code}), TreemapMix-PL
(Figure~\ref{fig:treemapmix-pl-code}), and the mixed-loss variant
(Figure~\ref{fig:treemapmix-mixed-code}) used in our experiments.

\begin{figure}[!htbp]
\centering
\footnotesize
\begin{verbatim}
import torch.nn.functional as F


def treemapmix_sce_loss(logits, labels, weights):
    """
    logits: model logits, shape (B, C)
    labels: target class indices, shape (B, K)
    weights: target area weights, shape (B, K)
    Each pair (labels[b, i], weights[b, i]) represents one
    TreemapMix region for sample b.
    """
    log_probs = F.log_softmax(logits, dim=1)
    target_log_probs = log_probs.gather(dim=1, index=labels)
    losses = -(weights * target_log_probs).sum(dim=1)
    return losses.mean()
\end{verbatim}
\caption{TreemapMix-SCE. The loss matches the area-weighted soft target induced by the TreemapMix regions.}
\label{fig:treemapmix-sce-code}
\end{figure}

\begin{figure}[t]
\centering
\footnotesize
\begin{verbatim}
import torch.nn.functional as F


def treemapmix_pl_loss(logits, labels, weights, eps=1e-12):
    """
    logits: model logits, shape (B, C)
    labels: ranked target class indices, shape (B, K)
    weights: ranked target area weights, shape (B, K)
    eps: small constant for numerical stability
    labels and weights are ordered from weakest to strongest
    target class.
    """
    log_probs = F.log_softmax(logits, dim=1)
    ranked_log_probs = log_probs.gather(dim=1, index=labels)
    ranked_probs = ranked_log_probs.exp()

    unranked_mass = 1.0 - ranked_probs.sum(dim=1, keepdim=True)
    denom_log = torch.log(unranked_mass + torch.cumsum(ranked_probs, dim=1) + eps)

    per_rank_losses = denom_log - ranked_log_probs
    losses = (weights * per_rank_losses).sum(dim=1)
    return losses.mean()
\end{verbatim}
\caption{TreemapMix-PL. The loss matches the area-induced ranking of the TreemapMix target classes.}
\label{fig:treemapmix-pl-code}
\end{figure}

\begin{figure}[!htbp]
\centering
\footnotesize
\begin{verbatim}
def treemapmix_mixed_loss(logits, labels, weights, beta=0.5):
    sce = treemapmix_sce_loss(logits, labels, weights)
    pl = treemapmix_pl_loss(logits, labels, weights)
    return beta * sce + (1.0 - beta) * pl
\end{verbatim}
\caption{TreemapMix-mixed, using the functions in Figures~\ref{fig:treemapmix-sce-code} and~\ref{fig:treemapmix-pl-code}. The coefficient $\beta$ matches the notation in the main text.}
\label{fig:treemapmix-mixed-code}
\end{figure}


\subsection{Runtime overhead}
\label{app:runtime-overhead}
For the configurations in this benchmark, augmentations are applied in the data loader.
Table~\ref{tab:augmentation-loader-runtime} isolates data-loader throughput
and host-memory usage for six augmentation configurations. The benchmark
measures CPU-side loading and augmentation without model forward or backward
passes. Peak resident set size (RSS) and unique set size (USS) are summed over
the loader process tree. RSS can count shared pages more than once, while USS
measures private memory. Memory is reported in GiB.

Relative to the single-image recipe (1282.8 images/s), TreemapMix-SCE at
$K=4$ reaches 1039.4 images/s, a 19.0\% throughput reduction;
TreemapMix-PL at $K=6$ reaches 954.3 images/s, a 25.6\% reduction.
The corresponding peak USS increases from 17.76 GiB to 20.16 and 21.74 GiB.
Mixup+CutMix reaches 1223.4 images/s, a 4.6\% reduction.
The RICAP configuration reaches 364.8 images/s with 25.18 GiB peak USS.
TreemapMix therefore adds measurable host-side cost.

\begin{table*}[t]
  \caption{Augmentation-loader throughput and peak host memory: mean $\pm$ sample standard deviation over five runs, with batch size 128, eight workers, 20 warm-up batches, and 100 measured batches. Throughput reduction is $100\times(1-T/T_{\mathrm{single}})$ using mean throughput. RSS and USS are process-tree resident and unique memory. The recorded Mosaic label does not establish a separate RICAP timing result.}
  \label{tab:augmentation-loader-runtime}
  \centering
  \footnotesize
  \setlength{\tabcolsep}{4pt}
  \begin{tabular}{@{}lcccc@{}}
    \toprule
    Method & \makecell{Throughput\\(images/s)} & \makecell{Throughput\\reduction} & \makecell{Peak RSS\\(GiB)} & \makecell{Peak USS\\(GiB)} \\
    \midrule
    No Aug & $1670.4 \pm 30.2$ & $-30.2\%$ & $28.49 \pm 0.07$ & $17.82 \pm 0.06$ \\
    Single-Image Aug & $1282.8 \pm 35.1$ & $+0.0\%$ & $28.38 \pm 0.16$ & $17.76 \pm 0.17$ \\
    Mixup+CutMix & $1223.4 \pm 29.7$ & $+4.6\%$ & $28.46 \pm 0.12$ & $17.82 \pm 0.15$ \\
    TreemapMix-SCE ($K=4$) & $1039.4 \pm 37.0$ & $+19.0\%$ & $30.85 \pm 0.34$ & $20.16 \pm 0.28$ \\
    TreemapMix-PL ($K=6$) & $954.3 \pm 33.0$ & $+25.6\%$ & $32.41 \pm 0.52$ & $21.74 \pm 0.49$ \\
    Mosaic (recorded) & $364.8 \pm 5.9$ & $+71.6\%$ & $36.82 \pm 0.14$ & $25.18 \pm 0.13$ \\
    \bottomrule
  \end{tabular}
\end{table*}


The training configuration uses a prefetch factor of 2, which allows CPU
preparation to overlap GPU computation. The isolated loader benchmark does
not measure how much of this overhead is hidden during training, so it does
not establish an end-to-end per-epoch timing comparison.

\clearpage
\section{Additional Experiments}
\subsection{Full ImageNet-1K Results}
\label{app:in1k-confidence-results}

Tables~\ref{tab:in1k-full-results-acc} and
\ref{tab:in1k-full-results-ece} report the full ImageNet-1K
main-comparison Top-1 accuracy and ECE results, respectively.

\begin{table*}[!htbp]
  \caption{Full ImageNet-1K Top-1 accuracy (percent) results under the matched 110-epoch protocol. Values are mean $\pm$ one sample standard deviation over seeds 42, 43, and 44. Each method uses one run per seed. Bold marks the best mean within each model column.}
  \label{tab:in1k-full-results-acc}
  \centering
  \footnotesize
  \setlength{\tabcolsep}{4pt}
  \begin{tabular}{@{}lcccc@{}}
    \toprule
    Method & ViT-Wee $\uparrow$ & ViT-Little $\uparrow$ & ViT-Medium $\uparrow$ & ViT-Betwixt $\uparrow$ \\
    \midrule
    No Aug & $75.10 \pm 0.24$ & $76.05 \pm 0.14$ & $76.59 \pm 0.02$ & $77.18 \pm 0.14$ \\
    Single-Image Aug & $77.94 \pm 0.15$ & $79.43 \pm 0.22$ & $79.80 \pm 0.04$ & $80.28 \pm 0.16$ \\
    Mixup+CutMix & $76.20 \pm 0.16$ & $78.84 \pm 0.05$ & $80.04 \pm 0.20$ & $81.09 \pm 0.13$ \\
    Mixup & $76.22 \pm 0.12$ & $78.64 \pm 0.10$ & $79.80 \pm 0.11$ & $80.84 \pm 0.09$ \\
    CutMix & $76.13 \pm 0.07$ & $79.12 \pm 0.07$ & $80.23 \pm 0.09$ & $81.23 \pm 0.12$ \\
    RICAP & $77.05 \pm 0.15$ & $78.41 \pm 0.14$ & $78.74 \pm 0.07$ & $79.13 \pm 0.11$ \\
    FMix & $76.73 \pm 0.10$ & $79.40 \pm 0.06$ & $80.30 \pm 0.07$ & $81.12 \pm 0.10$ \\
    GridMix & $76.39 \pm 0.17$ & $78.72 \pm 0.05$ & $79.50 \pm 0.09$ & $80.33 \pm 0.18$ \\
    ResizeMix & $76.96 \pm 0.17$ & $79.51 \pm 0.15$ & $80.53 \pm 0.09$ & $81.44 \pm 0.13$ \\
    SaliencyMix & $76.93 \pm 0.11$ & $79.53 \pm 0.07$ & $80.35 \pm 0.12$ & $81.22 \pm 0.06$ \\
    SmoothMix & $76.22 \pm 0.22$ & $79.27 \pm 0.09$ & $80.26 \pm 0.08$ & $81.18 \pm 0.05$ \\
    TokenMix & $76.43 \pm 0.10$ & $79.19 \pm 0.13$ & $80.27 \pm 0.13$ & $81.18 \pm 0.10$ \\
    TLA & $76.81 \pm 0.19$ & $79.47 \pm 0.16$ & $80.46 \pm 0.03$ & $81.32 \pm 0.13$ \\
    TreemapMix-SCE & $77.48 \pm 0.03$ & $79.60 \pm 0.16$ & $80.55 \pm 0.17$ & $81.22 \pm 0.15$ \\
    TreemapMix-mixed & $77.65 \pm 0.20$ & $79.68 \pm 0.17$ & $80.75 \pm 0.09$ & $81.30 \pm 0.05$ \\
    TreemapMix-PL & $\mathbf{78.01 \pm 0.13}$ & $\mathbf{80.12 \pm 0.04}$ & $\mathbf{81.05 \pm 0.14}$ & $\mathbf{81.69 \pm 0.16}$ \\
    \bottomrule
  \end{tabular}
\end{table*}


\begin{table*}[t]
  \caption{Full ImageNet-1K ECE (15 bins, percentage points) results under the matched 110-epoch protocol. Values are mean $\pm$ one sample standard deviation over seeds 42, 43, and 44. Each method uses one run per seed. Bold marks the best mean within each model column.}
  \label{tab:in1k-full-results-ece}
  \centering
  \footnotesize
  \setlength{\tabcolsep}{4pt}
  \begin{tabular}{@{}lcccc@{}}
    \toprule
    Method & ViT-Wee $\downarrow$ & ViT-Little $\downarrow$ & ViT-Medium $\downarrow$ & ViT-Betwixt $\downarrow$ \\
    \midrule
    No Aug & $3.62 \pm 0.13$ & $4.53 \pm 0.36$ & $4.29 \pm 0.20$ & $4.42 \pm 0.02$ \\
    Single-Image Aug & $6.29 \pm 0.13$ & $5.19 \pm 0.12$ & $5.53 \pm 0.21$ & $5.17 \pm 0.40$ \\
    Mixup+CutMix & $9.05 \pm 0.47$ & $7.53 \pm 0.10$ & $8.28 \pm 0.51$ & $7.91 \pm 0.35$ \\
    Mixup & $8.69 \pm 0.20$ & $7.54 \pm 0.15$ & $7.92 \pm 0.20$ & $7.70 \pm 0.11$ \\
    CutMix & $7.32 \pm 0.84$ & $5.56 \pm 0.54$ & $7.26 \pm 0.14$ & $5.63 \pm 0.20$ \\
    RICAP & $3.11 \pm 0.02$ & $4.43 \pm 0.11$ & $5.96 \pm 0.43$ & $6.44 \pm 0.11$ \\
    FMix & $6.40 \pm 0.19$ & $5.01 \pm 0.10$ & $5.44 \pm 0.26$ & $5.12 \pm 0.26$ \\
    GridMix & $7.04 \pm 0.15$ & $5.72 \pm 0.17$ & $5.66 \pm 0.54$ & $5.48 \pm 0.32$ \\
    ResizeMix & $7.18 \pm 0.28$ & $5.40 \pm 0.74$ & $6.06 \pm 0.17$ & $5.68 \pm 0.55$ \\
    SaliencyMix & $6.41 \pm 0.43$ & $5.09 \pm 0.09$ & $5.39 \pm 0.16$ & $4.57 \pm 0.29$ \\
    SmoothMix & $8.52 \pm 0.06$ & $6.93 \pm 0.18$ & $7.54 \pm 0.06$ & $6.88 \pm 0.21$ \\
    TokenMix & $7.60 \pm 0.65$ & $5.86 \pm 0.65$ & $5.89 \pm 0.75$ & $5.53 \pm 0.59$ \\
    TLA & $5.74 \pm 0.36$ & $4.68 \pm 0.03$ & $4.85 \pm 0.17$ & $4.53 \pm 0.36$ \\
    TreemapMix-SCE & $\mathbf{1.04 \pm 0.35}$ & $\mathbf{1.89 \pm 0.23}$ & $\mathbf{1.95 \pm 0.25}$ & $\mathbf{2.31 \pm 0.41}$ \\
    TreemapMix-mixed & $2.63 \pm 0.30$ & $3.54 \pm 0.05$ & $3.78 \pm 0.26$ & $4.26 \pm 0.13$ \\
    TreemapMix-PL & $5.63 \pm 0.18$ & $6.35 \pm 0.10$ & $6.45 \pm 0.05$ & $6.96 \pm 0.25$ \\
    \bottomrule
  \end{tabular}
\end{table*}


\subsection{Long-Tailed Recognition}
\label{app:long-tail-results}

We evaluate ImageNet-LT~\citep{tang2020longtailed} using the same four model
scales and compare Mixup+CutMix, RICAP, TreemapMix-SCE, and TreemapMix-PL.
For TreemapMix, we use the natural-frequency schedule: it independently
shuffles each class bucket and takes one sample per nonempty class per row
until each bucket is exhausted, preserving the original class frequencies.

TreemapMix-PL achieves the highest Top-1 accuracy on all four backbones,
improving over Mixup+CutMix by 1.40, 1.02, 1.31, and 1.82 percentage points
for ViT-Wee, Little, Medium, and Betwixt, respectively. Calibration behaves
differently from balanced ImageNet-1K: Mixup+CutMix has the lowest ECE on
every backbone, and neither TreemapMix objective retains that advantage.
TreemapMix-SCE also falls below the baseline in accuracy. Thus, the accuracy
benefit of PL extends to this imbalanced setting, while probability
calibration remains sensitive to the class distribution.

\begin{table*}[t]
  \caption{ImageNet-LT accuracy and calibration for four backbones. Values are mean $\pm$ one sample standard deviation over seeds 42 and 43. ECE uses 15 bins. Bold marks the best mean for each model and metric.}
  \label{tab:long-tail-results}
  \centering
  \footnotesize
  \setlength{\tabcolsep}{4pt}
  \begin{tabular}{@{}lcccc@{}}
    \toprule
    Method & ViT-Wee & ViT-Little & ViT-Medium & ViT-Betwixt \\
    \midrule
    \multicolumn{5}{c}{Top-1 accuracy (\%) $\uparrow$} \\
    \midrule
    Mixup+CutMix & $41.72 \pm 0.08$ & $41.67 \pm 0.43$ & $40.95 \pm 0.07$ & $39.78 \pm 0.21$ \\
    RICAP & $35.70 \pm 0.15$ & $35.39 \pm 0.00$ & $34.35 \pm 0.02$ & $31.44 \pm 3.37$ \\
    TreemapMix-SCE & $41.44 \pm 0.40$ & $41.09 \pm 0.40$ & $39.48 \pm 0.13$ & $38.67 \pm 0.27$ \\
    TreemapMix-PL & $\mathbf{43.12 \pm 0.06}$ & $\mathbf{42.69 \pm 0.07}$ & $\mathbf{42.26 \pm 0.07}$ & $\mathbf{41.60 \pm 0.17}$ \\
    \midrule
    \multicolumn{5}{c}{ECE (percentage points) $\downarrow$} \\
    \midrule
    Mixup+CutMix & $\mathbf{0.58 \pm 0.14}$ & $\mathbf{0.83 \pm 0.09}$ & $\mathbf{2.14 \pm 0.61}$ & $\mathbf{3.31 \pm 0.96}$ \\
    RICAP & $28.50 \pm 0.26$ & $31.86 \pm 0.21$ & $33.12 \pm 0.32$ & $30.89 \pm 2.81$ \\
    TreemapMix-SCE & $10.65 \pm 0.04$ & $13.00 \pm 0.56$ & $11.85 \pm 0.33$ & $11.35 \pm 0.62$ \\
    TreemapMix-PL & $17.90 \pm 0.63$ & $25.63 \pm 0.38$ & $23.62 \pm 0.74$ & $25.45 \pm 0.79$ \\
    \bottomrule
  \end{tabular}
\end{table*}


\subsection{Cross-Architecture Results}
\label{app:architecture-results}
To test whether the ImageNet-1K results extend beyond the four primary
ViT/EVA02 variants, we evaluate ConvNeXtV2-Nano, DeiT-Tiny, and Swin-Tiny
under the matched 110-epoch protocol.
We compare Mixup+CutMix, RICAP, TreemapMix-SCE, and TreemapMix-PL, retaining
the TreemapMix settings selected for the primary experiments.
Table~\ref{tab:architecture-results} reports accuracy and ECE.

PL gives the highest mean accuracy for each architecture, improving over
Mixup+CutMix by 0.66, 2.95, and 0.95 percentage points for ConvNeXtV2-Nano,
DeiT-Tiny, and Swin-Tiny. SCE has the lowest ECE for all three, at 2.74,
1.32, and 2.52 percentage points, respectively. The PL--SCE accuracy gap is
small on DeiT-Tiny and Swin-Tiny relative to the observed seed variation;
we report the mean ordering without claiming statistical significance.

\begin{table*}[htbp]
  \caption{ImageNet-1K results on three additional architectures under the 110-epoch protocol. Values are mean $\pm$ one sample standard deviation over seeds 42, 43, and 44. Accuracy is reported in percent, and ECE in percentage points using 15 bins. Bold marks the best mean for each model and metric.}
  \label{tab:architecture-results}
  \centering
  \footnotesize
  \setlength{\tabcolsep}{4pt}
  \begin{tabular}{@{}lcccccc@{}}
    \toprule
     & \multicolumn{2}{c}{ConvNeXtV2-Nano} & \multicolumn{2}{c}{DeiT-Tiny} & \multicolumn{2}{c}{Swin-Tiny} \\
    \cmidrule(lr){2-3} \cmidrule(lr){4-5} \cmidrule(lr){6-7}
    Method & Acc. $\uparrow$ & ECE $\downarrow$ & Acc. $\uparrow$ & ECE $\downarrow$ & Acc. $\uparrow$ & ECE $\downarrow$ \\
    \midrule
    Mixup+CutMix & $79.79 \pm 0.11$ & $9.48 \pm 0.13$ & $68.86 \pm 0.28$ & $12.34 \pm 0.27$ & $79.62 \pm 0.06$ & $7.50 \pm 0.13$ \\
    RICAP & $79.08 \pm 0.10$ & $5.95 \pm 0.17$ & $71.71 \pm 0.04$ & $2.04 \pm 0.28$ & $78.68 \pm 0.07$ & $7.78 \pm 0.05$ \\
    TreemapMix-SCE & $80.01 \pm 0.09$ & $\mathbf{2.74 \pm 0.27}$ & $71.78 \pm 0.26$ & $\mathbf{1.32 \pm 0.26}$ & $80.55 \pm 0.09$ & $\mathbf{2.52 \pm 0.11}$ \\
    TreemapMix-PL & $\mathbf{80.45 \pm 0.10}$ & $6.82 \pm 0.07$ & $\mathbf{71.81 \pm 0.07}$ & $2.71 \pm 0.15$ & $\mathbf{80.57 \pm 0.12}$ & $6.58 \pm 0.06$ \\
    \bottomrule
  \end{tabular}
\end{table*}

\subsection{CIFAR-100 Results}
\label{app:cifar100-results}
We evaluate all 16 augmentation settings on CIFAR-100 with the same four
ViT backbones, 256-pixel inputs, and the 110-epoch protocol. TreemapMix uses the settings from the primary
comparison: $K=4$ for SCE and the mixed objective, $K=6$ for PL, and a
cosine $\alpha$ schedule from $0.1$ to $0.5$.
Figure~\ref{fig:cifar100-results} reports Top-1 accuracy and ECE with 15 bins.

TreemapMix-SCE has the lowest ECE on all four backbones, ranging from
1.62 to 1.87 percentage points. TreemapMix-PL has the highest Top-1 accuracy
on ViT-Wee, Medium, and Betwixt, reaching 75.72\%, 75.97\%, and 76.39\%,
respectively. On ViT-Little, ResizeMix reaches 75.54\% compared with
75.42\% for PL. The mixed objective does not consistently improve on either
individual objective. These results retain the accuracy--calibration
trade-off, without a uniform accuracy advantage across every model.

\begin{figure*}[htbp]
  \centering
  \includegraphics{../evaluation/data/processed/figures/per_experiment/cifar100_per_experiment_short.pdf}
  \caption{CIFAR-100 comparison at 256-pixel resolution under the 110-epoch protocol. Points show the mean over seeds 42, 43, and 44; intervals denote $\pm$ one sample standard deviation. ECE uses 15 bins and is reported in percentage points. Average ranks summarize within-backbone ranks, with rank 1 best. SCE gives the lowest ECE throughout; PL gives the highest accuracy on three backbones, while ResizeMix leads on ViT-Little.}
  \label{fig:cifar100-results}
\end{figure*}


\subsection{Optimizer Robustness}
\label{app:optimizer-robustness}
We next test whether the gains of TreemapMix persist across different
optimizers on ImageNet-1K. We compare Muon
\citep{jordan2024muon}, AdamW (Adam with decoupled weight decay)
\citep{loshchilovDecoupledWeightDecay2019}, and NAdamW
\citep{Dahl2023AlgoPerf,Kasimbeg2025AlgoPerfResults}. We keep the core
TreemapMix hyperparameters fixed across optimizers. Because the main
experiments use NAdamW with a learning rate of $1.5\times10^{-3}$, we retain
this value for NAdamW and also use it for AdamW to keep the Adam-style
optimizer comparison matched. For Muon, we use a separately tuned learning
rate of $5\times10^{-3}$.

Table~\ref{tab:in1k-optimizer-results} compares Mixup+CutMix,
TreemapMix-SCE, and TreemapMix-PL under AdamW, Muon, and NAdamW for
the settings in the caption. These results should be interpreted as an
optimizer-sensitivity check rather than a definitive optimizer ranking. Across
the tested optimizers, the main qualitative pattern from the primary
ImageNet-1K comparison is preserved: TreemapMix-PL generally gives the highest
Top-1 accuracy, whereas TreemapMix-SCE gives the lowest ECE in most settings.
Muon gives the strongest Top-1 accuracy in several configurations, but this
gain is often accompanied by higher ECE. AdamW and NAdamW are more competitive
on calibration, although neither dominates across both metrics. Overall, these
single-seed results suggest that TreemapMix's accuracy and calibration trends
are not specific to one optimizer, while the absolute accuracy-calibration
trade-off still depends on the optimizer.

\begin{table*}[t]
  \centering
\caption{ImageNet-1K optimizer comparison under matched settings. Across AdamW,
Muon, and NAdamW, the same pattern largely holds: TreemapMix-PL tends to give
the highest accuracy, while TreemapMix-SCE gives the lowest ECE for every
model--optimizer pair. Accuracy is reported in percent and ECE in percentage
points; bold marks the best mean within each row.}
  \label{tab:in1k-optimizer-results}
  \begin{tabular}{llcccccc}
    \toprule
    & & \multicolumn{2}{c}{Mixup+CutMix}
      & \multicolumn{2}{c}{TreemapMix-SCE}
      & \multicolumn{2}{c}{TreemapMix-PL} \\
    \cmidrule(lr){3-4} \cmidrule(lr){5-6} \cmidrule(lr){7-8}
    Model & Optimizer
      & Acc. $\uparrow$ & ECE $\downarrow$
      & Acc. $\uparrow$ & ECE $\downarrow$
      & Acc. $\uparrow$ & ECE $\downarrow$ \\
    \midrule
    ViT-Wee
      & AdamW  & 75.76 & 8.99 & 77.18 & \textbf{0.43} & \textbf{77.75} & 5.41 \\
      & Muon   & 77.92 & 9.88 & 78.80 & \textbf{2.04} & \textbf{79.05} & 5.93 \\
      & NAdamW & 76.30 & 9.29 & 77.48 & \textbf{0.94} & \textbf{78.16} & 5.44 \\
    \midrule
    ViT-Little
      & AdamW  & 78.77 & 8.11 & 79.50 & \textbf{1.36} & \textbf{79.93} & 5.96 \\
      & Muon   & 80.27 & 8.61 & 80.72 & \textbf{3.05} & \textbf{80.87} & 6.57 \\
      & NAdamW & 78.80 & 7.65 & 79.72 & \textbf{1.90} & \textbf{80.09} & 6.23 \\
    \midrule
    ViT-Medium
      & AdamW  & 79.89 & 8.25 & 80.33 & \textbf{0.87} & \textbf{80.71} & 5.93 \\
      & Muon   & 80.98 & 7.86 & 81.42 & \textbf{2.77} & \textbf{81.43} & 6.49 \\
      & NAdamW & 80.04 & 7.69 & 80.73 & \textbf{2.09} & \textbf{81.01} & 6.39 \\
    \midrule
    ViT-Betwixt
      & AdamW  & 81.19 & 8.11 & 81.00 & \textbf{2.02} & \textbf{81.51} & 7.05 \\
      & Muon   & 81.71 & 7.32 & \textbf{81.85} & \textbf{3.32} & 81.79 & 7.15 \\
      & NAdamW & 81.13 & 7.50 & 81.26 & \textbf{2.75} & \textbf{81.80} & 7.18 \\
    \bottomrule
  \end{tabular}
\end{table*}

\subsection{ImageNet-1K Hyperparameter Sensitivity}
\label{app:in1k-hparam-sensitivity}

TreemapMix introduces two main hyperparameters: the Dirichlet concentration
$\alpha$ and the number of mixed regions $K$. We study their effects with a
single-seed hyperparameter sweep. The sweep identifies a stable operating
regime for both TreemapMix losses: smaller $\alpha$ values and small-to-medium
values of $K$ perform best, while more extreme compositions degrade both
accuracy and calibration. Overall, TreemapMix is more sensitive to $\alpha$
than to $K$.

We sweep
$\alpha \in \{0.2, 0.4, 0.6, 0.8, 1.0, 1.5, 2.0, 3.0\}$ and
$K \in \{2,\ldots,10\}$ using the same training hyperparameters as in the
main run. This sweep uses a simplified fixed-hyperparameter recipe: it does
not apply aspect-ratio filtering to TreemapMix regions and does not schedule
either $\alpha$ or $K$. Therefore, these results should be interpreted as a
robustness and sensitivity analysis rather than as a best-case configuration
search. The absolute numbers are consequently lower than the tuned main
results.

Figure~\ref{fig:in1k-hparam-sensitivity} shows the resulting sensitivity plots
and includes reference results for the main Mixup+CutMix and single-image-only
baselines. We include the single-image-only baseline because it is the
strongest non-TreemapMix setting in the main comparison for ViT-Wee.

\begin{figure*}[t]
  \centering
  \includegraphics{../evaluation/data/processed/figures/hyperparameter/hp_sensitivity_vit_wee_legend.pdf}\\[0pt]
  \includegraphics{../evaluation/data/processed/figures/hyperparameter/hp_sensitivity_vit_wee.pdf}\\[8pt]
  \includegraphics{../evaluation/data/processed/figures/hyperparameter/hp_sensitivity_k_vit_wee.pdf}
\caption{ImageNet-1K hyperparameter sensitivity of TreemapMix on ViT-Wee.
The top row sweeps $\alpha$ and marginalizes over $k$, while the bottom row
sweeps $k$ and marginalizes over $\alpha$. Curves show the median across the
marginalized hyperparameter, and shaded regions show the corresponding min--max
range. The best accuracy--calibration trade-off occurs for small $k$ and
moderately small $\alpha$. Full $(k,\alpha)$ heatmaps for Top-1 accuracy and
ECE are shown in Appendix Figure~\ref{fig:in1k-hparam-heatmaps-appendix}. These
experiments use the simplified sweep protocol without aspect-ratio filtering or
$k$/$\alpha$ scheduling.}
  \label{fig:in1k-hparam-sensitivity}
\end{figure*}


Figure~\ref{fig:in1k-hparam-heatmaps-appendix} reports the full
two-dimensional $(K,\alpha)$ heatmaps for the same sweep. These heatmaps show
the interaction between the number of regions and the Dirichlet concentration,
whereas Figure~\ref{fig:in1k-hparam-sensitivity} summarizes the marginal
trends.

\begin{figure*}[t]
  \centering
  \includegraphics{../evaluation/data/processed/figures/hyperparameter/hp_sensitivity_heatmaps_top1_vit_wee.pdf}\\[8pt]
  \includegraphics{../evaluation/data/processed/figures/hyperparameter/hp_sensitivity_heatmaps_ece_vit_wee.pdf}
  \caption{Appendix heatmaps for the ImageNet-1K TreemapMix hyperparameter
  sweep on ViT-Wee. ECE is measured using 15 bins. The top row shows Top-1 over the full $(K,\alpha)$ grid,
  and the bottom row shows ECE over the same grid. Both plots use the same
  simplified sweep protocol as the line plots in Figure~\ref{fig:in1k-hparam-sensitivity}: we do not apply
  aspect filtering and we do not use $\alpha$ or $K$ scheduling, so these
  absolute values are expected to be somewhat worse than the tuned main
  results. The heatmaps show that for best results in terms of ECE and accuracy, small $K$ values together with moderately small $\alpha$ values should be used.
  }
  \label{fig:in1k-hparam-heatmaps-appendix}
\end{figure*}


For Top-1 accuracy, both TreemapMix-SCE and TreemapMix-PL perform best at
relatively small $\alpha$ values and around $K=4$. Increasing $\alpha$ up to
approximately $1$ causes only modest degradation, while larger values lead to a
clearer accuracy drop. This trend suggests that TreemapMix benefits from some
imbalance among region areas rather than from nearly uniform region sizes.
Larger $K$ values also tend to reduce Top-1 accuracy, but this effect is
weaker and less uniform than the effect of $\alpha$. As $K$ increases,
performance varies more strongly across $\alpha$ values, indicating that
high-composition settings are more sensitive to the choice of mixing strength.

Calibration, measured by ECE with 15 bins, follows a similar but sharper
pattern. For both losses, ECE remains relatively stable up to a loss-dependent
threshold and then increases substantially. For TreemapMix-SCE, this threshold
occurs around $\alpha=1$; for TreemapMix-PL, it occurs around $\alpha=1.5$.
Increasing $K$ also tends to worsen calibration, but the effect is weaker and
less systematic than the effect of $\alpha$. Across the sweep, TreemapMix-PL
is less sensitive to hyperparameter variation, while TreemapMix-SCE achieves
the strongest calibration when $\alpha$ and $K$ are chosen carefully.

These results indicate that choosing $\alpha$ is more important than choosing
$K$ for reliable TreemapMix performance. For calibration-focused use, the
sweep supports TreemapMix-SCE with $\alpha \le 1$ and $K \le 5$. For
accuracy-focused use, the sweep supports TreemapMix-PL with $\alpha \le 0.4$
and $K \le 6$. The main recipe filters layouts with extreme patch aspect ratios and schedules
$\alpha$ from $0.1$ to $0.5$. These settings were not isolated in this sweep.

\subsection{Qualitative Visualizations and Attribution Analysis}
\label{app:qualitative-visualizations}

Figures~\ref{fig:treemapmix-randomness}, \ref{fig:treemapmix-k-sweep}, and
\ref{fig:treemapmix-alpha-sweep} visualize how TreemapMix changes with its
random draw and main hyperparameters. Figure~\ref{fig:treemapmix-randomness}
holds the source images and hyperparameters fixed while varying only the
random draw. Figures~\ref{fig:treemapmix-k-sweep} and
\ref{fig:treemapmix-alpha-sweep} sweep the number of regions and the
Dirichlet concentration, respectively, showing how these controls change the
granularity and balance of the mixed image.

\begin{figure*}[t]
  \centering
  \setlength{\tabcolsep}{2pt}
  \begin{tabular}{cccc}
    \includegraphics[width=0.235\textwidth]{../evaluation/data/processed/figures/augmentation_showcase/treemapmix_randomness/panel01_1024.pdf} &
    \includegraphics[width=0.235\textwidth]{../evaluation/data/processed/figures/augmentation_showcase/treemapmix_randomness/panel02_1024.pdf} &
    \includegraphics[width=0.235\textwidth]{../evaluation/data/processed/figures/augmentation_showcase/treemapmix_randomness/panel03_1024.pdf} &
    \includegraphics[width=0.235\textwidth]{../evaluation/data/processed/figures/augmentation_showcase/treemapmix_randomness/panel04_1024.pdf} \\
    \includegraphics[width=0.235\textwidth]{../evaluation/data/processed/figures/augmentation_showcase/treemapmix_randomness/panel05_1024.pdf} &
    \includegraphics[width=0.235\textwidth]{../evaluation/data/processed/figures/augmentation_showcase/treemapmix_randomness/panel06_1024.pdf} &
    \includegraphics[width=0.235\textwidth]{../evaluation/data/processed/figures/augmentation_showcase/treemapmix_randomness/panel07_1024.pdf} &
    \includegraphics[width=0.235\textwidth]{../evaluation/data/processed/figures/augmentation_showcase/treemapmix_randomness/panel08_1024.pdf}
  \end{tabular}
\caption{TreemapMix randomness with fixed source images and hyperparameters
($\alpha=0.5$, $K=5$). Each panel uses a different random seed for the
Dirichlet draw, showing that TreemapMix produces diverse spatial compositions
while preserving the same sources and mixing settings.}
  \label{fig:treemapmix-randomness}
\end{figure*}


\begin{figure*}[t]
  \centering
  \setlength{\tabcolsep}{2pt}
  \begin{tabular}{cccc}
    \includegraphics{../evaluation/data/processed/figures/augmentation_showcase/treemapmix_k_sweep/k02_1024.pdf} &
    \includegraphics{../evaluation/data/processed/figures/augmentation_showcase/treemapmix_k_sweep/k03_1024.pdf} &
    \includegraphics{../evaluation/data/processed/figures/augmentation_showcase/treemapmix_k_sweep/k04_1024.pdf} &
    \includegraphics{../evaluation/data/processed/figures/augmentation_showcase/treemapmix_k_sweep/k05_1024.pdf} \\
    $K=2$ & $K=3$ & $K=4$ & $K=5$ \\
    \includegraphics{../evaluation/data/processed/figures/augmentation_showcase/treemapmix_k_sweep/k06_1024.pdf} &
    \includegraphics{../evaluation/data/processed/figures/augmentation_showcase/treemapmix_k_sweep/k07_1024.pdf} &
    \includegraphics{../evaluation/data/processed/figures/augmentation_showcase/treemapmix_k_sweep/k08_1024.pdf} &
    \includegraphics{../evaluation/data/processed/figures/augmentation_showcase/treemapmix_k_sweep/k09_1024.pdf} \\
    $K=6$ & $K=7$ & $K=8$ & $K=9$
  \end{tabular}
  \caption{Effect of the number of TreemapMix source regions. All panels use the same source images with fixed TreemapMix hyperparameters $\alpha=2.0$ and maximum aspect ratio $15.0$; increasing $K$ from $2$ to $9$ produces finer-grained spatial mixtures.}
  \label{fig:treemapmix-k-sweep}
\end{figure*}


\begin{figure*}[t]
  \centering
  \setlength{\tabcolsep}{2pt}
  \begin{tabular}{cccc}
    \includegraphics[width=0.235\textwidth]{../evaluation/data/processed/figures/augmentation_showcase/treemapmix_alpha_sweep/alpha0_05_1024.pdf} &
    \includegraphics[width=0.235\textwidth]{../evaluation/data/processed/figures/augmentation_showcase/treemapmix_alpha_sweep/alpha0_1_1024.pdf} &
    \includegraphics[width=0.235\textwidth]{../evaluation/data/processed/figures/augmentation_showcase/treemapmix_alpha_sweep/alpha0_3_1024.pdf} &
    \includegraphics[width=0.235\textwidth]{../evaluation/data/processed/figures/augmentation_showcase/treemapmix_alpha_sweep/alpha0_5_1024.pdf} \\
    $\alpha=0.05$ & $\alpha=0.1$ & $\alpha=0.3$ & $\alpha=0.5$ \\[2pt]
    \includegraphics[width=0.235\textwidth]{../evaluation/data/processed/figures/augmentation_showcase/treemapmix_alpha_sweep/alpha1_1024.pdf} &
    \includegraphics[width=0.235\textwidth]{../evaluation/data/processed/figures/augmentation_showcase/treemapmix_alpha_sweep/alpha1_5_1024.pdf} &
    \includegraphics[width=0.235\textwidth]{../evaluation/data/processed/figures/augmentation_showcase/treemapmix_alpha_sweep/alpha3_1024.pdf} &
    \includegraphics[width=0.235\textwidth]{../evaluation/data/processed/figures/augmentation_showcase/treemapmix_alpha_sweep/alpha5_1024.pdf} \\
    $\alpha=1.0$ & $\alpha=1.5$ & $\alpha=3.0$ & $\alpha=5.0$ \\[2pt]
  \end{tabular}
  \caption{Effect of the Dirichlet concentration in TreemapMix. All panels use the same source images with fixed TreemapMix hyperparameters $K=5$ and maximum aspect ratio $15.0$; smaller $\alpha$ values increase the likelihood of images being distributed unequally. }
  \label{fig:treemapmix-alpha-sweep}
\end{figure*}


We also complement the quantitative logit-alignment diagnostics with
qualitative attribution maps in Figure~\ref{fig:feature-space-analysis}. The
maps are computed on the same held-out diagnostic set using the ViT-Betwixt
checkpoints. We use Gradient$\times$Input attribution
\citep{shrikumarNotJustBlack2017a,simonyanDeepConvolutionalNetworks2014},
which multiplies the input by the gradient of the target logit with respect to
the input. This provides a qualitative view of which source regions contribute
most strongly to a given class logit.

Figure~\ref{fig:feature-space-analysis} highlights a qualitative difference
between the two TreemapMix losses. For the boathouse example, TreemapMix-SCE assigns a displayed probability
of 0.01, while the other methods round to 0.00; rounded zeros do not imply
exactly zero probability. By contrast,
TreemapMix-PL tends to assign substantially higher probability to the largest
region. This is consistent with the ranking objective, which emphasizes relative
ordering rather than matching the target probabilities.

The attribution overlays are often more focused on the target regions for the
TreemapMix-trained models, including the suit and boathouse examples with smaller target areas. In contrast, the non-TreemapMix baselines often concentrate
on more localized discriminative features. For example, in the Cougar row, the
baseline attributions focus more strongly on the face. These qualitative
results are not intended as standalone evidence, but they are consistent with
the quantitative logit-alignment diagnostics: TreemapMix changes how the model
distributes evidence across composed regions, and the PL objective makes this
effect most explicit.

\begin{figure*}[t]
  \centering
  \includegraphics[width=\textwidth]{../evaluation/data/processed/figures/diagnostic_attribution/class_286/figure.pdf}\\[3pt]
  \includegraphics[width=\textwidth]{../evaluation/data/processed/figures/diagnostic_attribution/class_834/figure.pdf}\\[3pt]
  \includegraphics[width=\textwidth]{../evaluation/data/processed/figures/diagnostic_attribution/class_449/figure.pdf}
\caption{Diagnostic attribution visualizations for three ImageNet-1K target
classes. Each row shows the boxed input followed by Gradient$\times$Input
overlays for Mixup+CutMix, RICAP,
TreemapMix-SCE, and TreemapMix-PL. Red boxes mark the target-class source patch;
brighter regions indicate larger attribution magnitude for the corresponding
target-class logit. Values below panels show predicted target-class
probabilities, and row labels report target patch area. In these examples,
TreemapMix variants more often concentrate attribution on the target patch;
TreemapMix-PL produces sharper probabilities, whereas TreemapMix-SCE remains
more moderate.}
  \label{fig:feature-space-analysis}
\end{figure*}

